%
% einleitung.tex -- Beispiel-File für die Einleitung
%
% (c) 2020 Prof Dr Andreas Müller, Hochschule Rapperswil
%
% !TEX root = ../../buch.tex
% !TEX encoding = UTF-8
%
\section{Motivation und Definition
\label{hauptwert:section:motivation-definition}}
\kopfrechts{Teil 0}
\subsection{Uneigentliche Integrale}

Man betrachte die Funktion
\begin{equation}
f(x) = \frac{1}{x}
\end{equation}
und damit das Integral
\begin{equation}
\int_{-\infty}^{\infty} f(x)\,dx
= \int_{-\infty}^{\infty} \frac{1}{x}\,dx.
\end{equation}
%
% kehrwert.tex
%
% (c) 2026 Damien Flury, OST Ostschweizer Fachhochschule
%
\begin{figure}
\centering
\includegraphics{papers/hauptwert/images/kehrwert.pdf}
\caption{Graph der Funktion $f(x)=1/x$.
Die Funktion hat bei $x=0$ eine Polstelle:
für $x\to 0^+$ wächst sie unbeschränkt,
für $x\to 0^-$ fällt sie unbeschränkt.
Die beiden Äste sind punktsymmetrisch zum Ursprung.
\label{hauptwert:figure:kehrwert}}
\end{figure}

Abbildung \ref{hauptwert:figure:kehrwert} zeigt den Graphen des Integranden
$f(x)=1/x$.
Beim Lösen dieses Integrals fällt schnell auf, dass es ein Problem gibt, und
zwar die Polstelle. Die Funktion hat eine nicht-hebbare
Definitionslücke bei $x = 0$ und beide Seiten divergieren.

\subsection{Ein erster Versuch}
Um das Integral zu lösen, wäre der erste natürliche Schritt, beide Seiten
einzeln zu Lösen und mit dem Grenzwert sich der Polstelle zu nähern:

\begin{equation}
\int_{-\infty}^{\infty} f(x)\,dx
= \lim_{\epsilon_1 -> 0^-} \int_{-\infty}^{\epsilon_1} f(x)\,dx + \lim_{\epsilon_2 -> 0^+} \int_{\epsilon_2}^{\infty} f(x)\,dx
\end{equation}
